% ===== Block B00b: 原文 p.vii (Notation) =====
% [原文 p.vii]
% [Tutor Note: 这张表是全书唯一的符号总表，第 2、3 章几乎每条公式都要回来查。
%  建议精读前先抄一遍：hat g_N 的直和结构、kappa、c_v、Delta_lambda 四个量
%  决定了共形块的维数与射影联络的具体形式。]
\begingroup
\setlength{\parindent}{0pt}
\setlength{\parskip}{0.55em}
\newcommand{\nota}[3][4.0cm]{\hangindent=#1\hangafter=1\makebox[#1][l]{#2}#3\par}

\nota{$S_\omega$}{与对称双微分 $\omega$ 相伴的射影联络（见式 (1.41)），}
\nota{$T_X$}{复流形（或光滑射影簇）$X$ 的全纯（或正则）切丛，}
\nota{$\Theta_X$}{复流形或光滑射影簇 $X$ 的全纯向量场层，}
\nota{$\Omega_X^k$}{复流形或光滑射影簇 $X$ 的全纯 $k$-形式层，}
\nota{$\hat{\mathfrak g}=\mathfrak g\otimes_\C\C((\xi))\oplus\C\cdot c$,}{}
\noindent$\hat{\mathfrak g}_N=\mathfrak g\otimes_\C\Bigl(\bigoplus_{j=1}^{N}\C((\xi_j))\Bigr)\oplus\C\cdot c$,\par
\nota{$\Delta$}{复单李代数及其 Cartan 子代数 $(\mathfrak g,\mathfrak h)$ 的根系，}
\nota{$\Delta_+$}{正根之集，}
\nota{$\theta$}{复单李代数 $\mathfrak g$ 的最长根，}
\nota{$\rho=\frac12\sum_{\alpha\in\Delta_+}\alpha$,}{}
\nota{$(\ ,\ )$}{复单李代数 $\mathfrak g$ 的 Cartan--Killing 型，归一化使
  $(\theta,\theta)=2$，}
\nota{$P_+$}{复单李代数 $\mathfrak g$ 的整支配权之集，}
\nota{$P_\ell=\{\lambda\in P_+\mid(\lambda,\theta)\le\ell\}$,}{}
\nota{$\mu^\dagger=-w(\mu)$}{其中 $\mu\in P_\ell$，$w$ 为复单李代数 $\mathfrak g$ 的 Weyl 群的最长元，}
\nota{$\kappa=g^*+\ell$}{其中 $g^*$ 为复单李代数 $\mathfrak g$ 的对偶 Coxeter 数
  （当 $\mathfrak g=\mathfrak{sl}(n+1,\C)$ 时 $\kappa=n+1+\ell$），}
\nota{$c_v=\dfrac{\ell\dim\mathfrak g}{\kappa}$,}{}
\nota{$\Delta_\lambda=\dfrac{(\hat\lambda,\lambda+2\rho)}{2\kappa}$,}{}
\nota{$\hat\Delta_v=\Delta_\lambda+\Delta_{\mu_1}-\Delta_{\mu_2}$,}{}
\nota[7.0cm]{$\hat\Delta_{\mathbf v}=\Delta_\lambda+\Delta_{\mu_1}-\Delta_{\mu_2}$}{对一个顶点
  $\mathbf v=\binom{\lambda}{\mu_1\ \mu_2}$，}
\nota[6.6cm]{$X_n=\C^n\setminus\bigcup_{i<j}\Delta_{ij}$,\quad
  $\Delta_{ij}=\{(z_n,\dots,z_1)\in\C^n\mid z_i=z_j\}$,}{}
\nota[6.6cm]{$\mathcal R_n=\{(z_n,\dots,z_1)\in\C^n\mid |z_n|>\cdots>|z_1|\}$,}{}
\endgroup

% [Tutor Note: 原文此处把 $\hat\Delta_v$ 的定义列了两遍（前一行只给公式，后一行补
%  "for a vertex v = (lambda; mu_1, mu_2)" 的说明）。照原文保留，见 ERRATA。
%  $\mathcal R_n$ 是第 3 章"标准的"有序参数化区域的记号，$X_n$ 是去掉所有对角
%  因子的位形空间；$|z_n|>\cdots>|z_1|$ 是单值化（展开）讨论的起点。]
